\documentclass[../../main.tex]{subfiles}
\graphicspath{{../../image/}{../../figure/}}
\begin{document}

\section{Riemann 不变量}\label{section:ch07-Riemann不变量}

考虑 $\mathscr{O}\subset\mathbb{R}^n$ 上的双曲守恒律组\eqref{align:ch07-典范守恒律}。一个非常重要的概念由下述内容引入。

\begin{definition}[$i$-Riemann 不变量]\label{definition:ch07-iRiemann不变量}
\eqref{align:ch07-典范守恒律}的一个\emph{$i$-Riemann 不变量}是 $\mathscr{O}$ 上的光滑标量值函数 $w$，满足
\begin{align}
\mathrm{D}w(U)R_i(U)=0,\qquad U\in\mathscr{O}.
\label{align:ch07-Riemann不变量定义}
\end{align}
\end{definition}

例如，回忆\eqref{align:ch07-p系统右特征}，可直接验证函数
\begin{align}
w=-\int^u\sqrt{\sigma'(\omega)}\,\mathrm{d}\omega+v,\qquad
z=-\int^u\sqrt{\sigma'(\omega)}\,\mathrm{d}\omega-v
\label{align:ch07-p系统Riemann}
\end{align}
分别是方程组\eqref{align:ch07-p系统}的 1- 与 2-Riemann 不变量。类似地，可以证明
\begin{align}
w=v+\frac{2(\kappa\gamma)^{1/2}}{\gamma-1}\rho^{\frac{\gamma-1}{2}},\qquad
z=v-\frac{2(\kappa\gamma)^{1/2}}{\gamma-1}\rho^{\frac{\gamma-1}{2}}
\label{align:ch07-理想气体Riemann}
\end{align}
是理想气体等熵流动方程组\eqref{align:ch07-理想气体}的 1- 与 2-Riemann 不变量。\footnote{在等温情形 $\gamma=1$，$w=v+\kappa^{1/2}\log\rho$，$z=v-\kappa^{1/2}\log\rho$。}

通过求解关于 $w$ 的一阶线性微分方程\eqref{align:ch07-Riemann不变量定义}，可以在任一点 $U\in\mathscr{O}$ 的邻域内构造 $n-1$ 个 $i$-Riemann 不变量，其梯度线性无关并张成 $R_i$ 的正交补。例如，读者可作为习题验证下述三对函数
\begin{align}
\begin{cases}
s,\ -\displaystyle\int^u\sqrt{-p_\omega(\omega,s)}\,\mathrm{d}\omega+v\\[2mm]
v,\ -p(u,s)\\[2mm]
s,\ -\displaystyle\int^u\sqrt{-p_\omega(\omega,s)}\,\mathrm{d}\omega-v
\end{cases}
\label{align:ch07-热弹性Riemann}
\end{align}
分别是绝热热弹性方程组\eqref{align:ch07-热弹性}的 1-、2- 与 3-Riemann 不变量。

Riemann 不变量在具有如下特殊结构的方程组中特别有用。

\begin{definition}[Riemann 不变量坐标系]\label{definition:ch07-Riemann坐标系}
方程组\eqref{align:ch07-典范守恒律}\emph{配有一个 Riemann 不变量坐标系}，如果存在 $\mathscr{O}$ 上的 $n$ 个标量值函数 $(w_1,\cdots,w_n)$，使得对任意 $i,j=1,\cdots,n$，$i\neq j$，$w_j$ 是\eqref{align:ch07-典范守恒律}的一个 $i$-Riemann 不变量。
\end{definition}

\hyperref[definition:ch07-iRiemann不变量]{定义 7.3.1}与\hyperref[definition:ch07-Riemann坐标系]{定义 7.3.2}的一个直接推论是：

\begin{theorem}\label{theorem:ch07-Riemann坐标系等价}
函数 $(w_1,\cdots,w_n)$ 构成\eqref{align:ch07-典范守恒律}的 Riemann 不变量坐标系，当且仅当
\begin{align}
\mathrm{D}w_i(U)R_j(U)=
\begin{cases}
=0,& i\neq j\\
\neq 0,& i=j,
\end{cases}
\label{align:ch07-Riemann坐标系条件}
\end{align}
即，当且仅当对 $i=1,\cdots,n$，$\mathrm{D}w_i(U)$ 是矩阵 $\mathrm{D}F(U)$ 的关联于特征速度 $\lambda_i(U)$ 的左特征向量。等价地，$w_i$ 的等值面在任一点 $U$ 处的切超平面由向量 $R_1(U),\cdots,R_{i-1}(U),R_{i+1}(U),\cdots,R_n(U)$ 张成。
\end{theorem}

假定\eqref{align:ch07-典范守恒律}配有 Riemann 不变量坐标系 $(w_1,\cdots,w_n)$，从左乘 $\mathrm{D}w_i$（$i=1,\cdots,n$），我们把该方程组约化为对角形式：
\begin{align}
\partial_t w_i+\lambda_i\partial_x w_i=0,\qquad i=1,\cdots,n,
\label{align:ch07-对角化}
\end{align}
这等价于原形式\eqref{align:ch07-典范守恒律}，尽管仅在古典解的框架内。\eqref{align:ch07-对角化}的左端恰是 $w_i$ 在 $i$-特征方向上的导数。因而，

\begin{theorem}\label{theorem:ch07-Riemann不变量沿特征线常值}
假定 $(w_1,\cdots,w_n)$ 构成\eqref{align:ch07-典范守恒律}的 Riemann 不变量坐标系。对 $i=1,\cdots,n$，在\eqref{align:ch07-典范守恒律}的任一古典解 $U$ 所关联的每条 $i$-特征线上，$w_i$ 保持常值。
\end{theorem}

显然，任何两个守恒律的双曲方程组都配有 Riemann 不变量坐标系。相反，在规模 $n\geqslant 3$ 的方程组中，Riemann 不变量坐标系仅在如下例外情形中存在：形式上过定的方程组\eqref{align:ch07-Riemann坐标系条件}（关于 $n$ 个未知量 $(w_1,\cdots,w_n)$ 有 $n(n-1)$ 个方程）有解。由 Frobenius 定理，$w_i$ 等值面的切超平面由 $R_1,\cdots,R_{i-1},R_{i+1},\cdots,R_n$ 张成，当且仅当对 $i\neq j\neq k\neq i$，Lie 括号 $[R_j,R_k]$（见\eqref{align:ch07-Lie括号定义}）落在 $\{R_1,\cdots,R_{i-1},R_{i+1},\cdots,R_n\}$ 的张成中。因而，方程组\eqref{align:ch07-典范守恒律}配有 Riemann 不变量坐标系，当且仅当
\begin{align}
[R_j,R_k]=\alpha_j^k R_j-\alpha_k^j R_k,\qquad j,k=1,\cdots,n,
\label{align:ch07-Frobenius条件}
\end{align}
其中 $\alpha_i^\ell$ 是标量场。

当\eqref{align:ch07-典范守恒律}存在 Riemann 不变量坐标系 $(w_1,\cdots,w_n)$ 时，方便的做法是把特征向量 $R_1,\cdots,R_n$ 规范化，使得
\begin{align}
\mathrm{D}w_i(U)R_j(U)=
\begin{cases}
0,& i\neq j\\
1,& i=j.
\end{cases}
\label{align:ch07-Riemann规范化}
\end{align}
在此情形下，我们注意恒等式
\begin{align}
\mathrm{D}w_i\mathrm{D}R_j R_k=\mathrm{D}(\mathrm{D}w_i R_j)R_k-R_j^\top\mathrm{D}^2 w_i R_k=-R_j^\top\mathrm{D}^2 w_i R_k,\qquad i,j,k=1,\cdots,n,
\label{align:ch07-DwDR}
\end{align}
它特别蕴含 $\mathrm{D}w_i[R_j,R_k]=0$，$i=1,\cdots,n$，即
\begin{align}
[R_j,R_k]=0,\qquad j,k=1,\cdots,n.
\label{align:ch07-Lie括号零}
\end{align}

回忆恒等式\eqref{align:ch07-Lie括号恒等式}并利用\eqref{align:ch07-Lie括号定义}与\eqref{align:ch07-Lie括号零}，我们推出：只要 $\lambda_j(U)\neq\lambda_k(U)$，$\mathrm{D}R_j(U)R_k(U)$ 就落在 $\{R_j(U),R_k(U)\}$ 的张成中。这与\eqref{align:ch07-Riemann规范化}、\eqref{align:ch07-DwDR}一起给出
\begin{align}
R_j^\top\mathrm{D}^2 w_i R_k=-\mathrm{D}w_i\mathrm{D}R_j R_k=0,\qquad i\neq j\neq k\neq i.
\label{align:ch07-D2w交叉}
\end{align}

当\eqref{align:ch07-典范守恒律}拥有 Riemann 不变量坐标系 $(w_1,\cdots,w_n)$ 时，把 $U$ 映到 $W=(w_1,\cdots,w_n)^\top$ 的映射是局部微分同胚。把 $W$ 而非 $U$ 视为状态向量常常是方便的。为避免符号泛滥，当无混淆之虞时，我们将使用同一符号表示作为 $U$ 或 $W$ 函数的场。由\eqref{align:ch07-Riemann规范化}，$\partial U/\partial w_i=R_i$，于是对典型函数 $\phi$，链式法则给出
\begin{align}
\frac{\partial\phi}{\partial w_i}=\mathrm{D}\phi R_i,\qquad i=1,\cdots,n.
\label{align:ch07-链式法则}
\end{align}
例如，\eqref{align:ch07-Lie括号零}退化为 $\partial R_j/\partial w_k=\partial R_k/\partial w_j=\partial^2 U/\partial w_j\partial w_k$。

我们进而推导某些恒等式，它们稍后将帮助我们确立配有 Riemann 不变量坐标系的方程组的其他显著性质。合并\eqref{align:ch07-Lie括号恒等式}、\eqref{align:ch07-Lie括号定义}、\eqref{align:ch07-Lie括号零}与\eqref{align:ch07-链式法则}，我们推出
\begin{align}
-\frac{\partial R_j}{\partial w_k}=g_{jk}R_j+g_{kj}R_k,\qquad j,k=1,\cdots,n;\ j\neq k,
\label{align:ch07-Rj导数}
\end{align}
其中我们置
\begin{align}
g_{jk}=\frac{1}{\lambda_j-\lambda_k}\frac{\partial\lambda_j}{\partial w_k},\qquad j,k=1,\cdots,n;\ j\neq k.
\label{align:ch07-gjk定义}
\end{align}
注意即使当 $\lambda_j=\lambda_k$ 时 $g_{jk}$ 也可能有定义，因为在这类点上 $\partial\lambda_j/\partial w_k$ 由\eqref{align:ch07-退化点Dlambda}与\eqref{align:ch07-链式法则}而知为零。由\eqref{align:ch07-Rj导数}，
\begin{align}
-\frac{\partial^2 R_j}{\partial w_i\partial w_k}
=\frac{\partial g_{jk}}{\partial w_i}R_j-g_{jk}(g_{ji}R_j+g_{ij}R_i)
+\frac{\partial g_{kj}}{\partial w_i}R_k-g_{kj}(g_{ki}R_k+g_{ik}R_i).
\label{align:ch07-Rj二阶导}
\end{align}
由于对 $i\neq j\neq k\neq i$，$R_i,R_j,R_k$ 线性无关，且\eqref{align:ch07-Rj二阶导}的右端关于 $(i,k)$ 必须对称，我们推出
\begin{align}
\frac{\partial g_{jk}}{\partial w_i}=\frac{\partial g_{ji}}{\partial w_k},\qquad i\neq j\neq k\neq i,
\label{align:ch07-g对称}
\end{align}
\begin{align}
\frac{\partial g_{ij}}{\partial w_k}+g_{ij}g_{jk}-g_{ij}g_{ik}+g_{ik}g_{kj}=0,\qquad i\neq j\neq k\neq i.
\label{align:ch07-gCocycle}
\end{align}

在应用中出现的规模 $n\geqslant 3$ 的双曲守恒律组中，很少配有 Riemann 不变量坐标系。一个值得注意的例子是\emph{电泳（electrophoresis）}方程组：
\begin{align}
\partial_t U_i+\partial_x\left[\frac{c_i U_i}{\sum_{j=1}^n U_j}\right]=0,\qquad i=1,\cdots,n,
\label{align:ch07-电泳}
\end{align}
其中 $c_1<c_2<\cdots<c_n$ 是正常数。该方程组控制通过施加电场分离溶液中 $n$ 种离子化合物的过程。在此背景下，$U_i$ 表示浓度，$c_i$ 度量第 $i$ 种的电泳迁移率。特别地，$U_i\geqslant 0$。作为习题，读者可验证\eqref{align:ch07-电泳}的特征速度由
\begin{align}
\lambda_i=\mu_i\sum_{j=1}^n U_j,\qquad i=1,\cdots,n,
\label{align:ch07-电泳特征速度}
\end{align}
给出，其中对 $i=1,\cdots,n-1$，$U$ 处 $\mu_i$ 的值是方程
\begin{align}
\sum_{j=1}^n\frac{c_j U_j}{c_j-\mu}=\sum_{j=1}^n U_j
\label{align:ch07-电泳mu方程}
\end{align}
在区间 $(c_i,c_{i+1})$ 中的解；而 $\mu_n=0$。进而，\eqref{align:ch07-电泳}配有 Riemann 不变量坐标系 $(w_1,\cdots,w_n)$，其中对 $i=1,\cdots,n-1$，$U$ 处 $w_i$ 的值是方程
\begin{align}
\sum_{j=1}^n\frac{U_j}{c_j-w}=0
\label{align:ch07-电泳wi方程}
\end{align}
在区间 $(c_i,c_{i+1})$ 中的解；以及
\begin{align}
w_n=\sum_{j=1}^n\frac{1}{c_j}U_j.
\label{align:ch07-电泳wn}
\end{align}
稍后我们将看到，方程组\eqref{align:ch07-电泳}具有非常特殊的结构与一系列有趣性质。

另一个配有 Riemann 不变量坐标系的有趣方程组是\eqref{align:ch07-电磁波}，我们回忆它控制平面电磁波通过特殊各向同性电介质的传播。这通过从 $(B_1,B_2,D_1,D_2)$ 过渡到由
\begin{align}
\begin{cases}
\sqrt{2}p\exp(\mathrm{i}a)=B_2+D_1-\mathrm{i}(B_1-D_2)\\
\sqrt{2}q\exp(\mathrm{i}b)=-B_2+D_1+\mathrm{i}(B_1+D_2)
\end{cases}
\label{align:ch07-电磁复变换}
\end{align}
定义的新状态向量 $(p,q,a,b)$ 而看出。特别地，$p^2+q^2=r^2$。简单计算表明，至少在古典解的框架内，\eqref{align:ch07-电磁波}退化为
\begin{align}
\begin{cases}
\partial_t p+\partial_x\left[\dfrac{\psi'(r)}{r}p\right]=0\\
\partial_t q-\partial_x\left[\dfrac{\psi'(r)}{r}q\right]=0,
\end{cases}
\label{align:ch07-电磁pq}
\end{align}
\begin{align}
\begin{cases}
\partial_t a+\dfrac{\psi'(r)}{r}\partial_x a=0\\
\partial_t b-\dfrac{\psi'(r)}{r}\partial_x b=0.
\end{cases}
\label{align:ch07-电磁ab}
\end{align}

注意\eqref{align:ch07-电磁pq}构成两个守恒律的封闭方程组，由此可确定 $p,q$ 从而 $r$。继而\eqref{align:ch07-电磁ab}可作为两个独立的非齐次标量守恒律求解，以确定 $a$ 与 $b$。特别地，$a,b$ 连同\eqref{align:ch07-电磁pq}的任一对 Riemann 不变量，将构成\eqref{align:ch07-电磁波}的一个 Riemann 不变量坐标系。

\end{document}
